% ============================================================
%  C: Como Programar -- 6ª Edição | Deitel & Deitel
%  Capítulo 13 -- O Pré-processador em C
%  EXERCÍCIOS DE AUTORREVISÃO (13.1 -- 13.3)
%  Abordagem incremental: Caps. 1 a 13
% ============================================================
\documentclass[12pt,a4paper]{article}
\usepackage[utf8]{inputenc}
\usepackage[T1]{fontenc}
\usepackage[brazil]{babel}
\usepackage{geometry}
\geometry{top=2.5cm,bottom=2.5cm,left=3cm,right=3cm}
\usepackage{parskip}\setlength{\parskip}{6pt}\setlength{\parindent}{0pt}
\usepackage{xcolor,tcolorbox,enumitem,listings,fancyhdr,titlesec,hyperref,booktabs,array}
\tcbuselibrary{skins,breakable}

\definecolor{deitelblue}{RGB}{0,70,127}
\definecolor{deitelgray}{RGB}{240,242,245}
\definecolor{deitelgreen}{RGB}{0,110,60}
\definecolor{deitellightblue}{RGB}{220,235,250}
\definecolor{deitelred}{RGB}{160,20,20}
\definecolor{deitellorange}{RGB}{200,90,0}

\newtcolorbox{boaprática}[1][]{colback=deitellightblue,colframe=deitelblue,
  fonttitle=\bfseries\small,title={Boa prática de programação},breakable,#1}
\newtcolorbox{erroco}[1][]{colback=red!5,colframe=deitelred,
  fonttitle=\bfseries\small,title={Erro comum de programação},breakable,#1}
\newtcolorbox{respostabox}[1][]{colback=deitelgray,colframe=deitelblue!60,
  fonttitle=\bfseries\small\color{deitelblue},title={Resposta detalhada},breakable,#1}
\newtcolorbox{enunciadoBox}[1][]{colback=white,colframe=deitelblue,
  fonttitle=\bfseries,title={#1},breakable}
\newtcolorbox{saida}[1][]{colback=black!88,colframe=deitelblue,
  fontupper=\ttfamily\small\color{green!70},title={\small\color{white}Saída do programa},breakable,#1}

\setlist[enumerate,1]{label=\textbf{\alph*)},leftmargin=2em}
\setlist[itemize]{leftmargin=2em}

\lstset{
  basicstyle=\ttfamily\small,backgroundcolor=\color{deitelgray},
  frame=single,framerule=0.4pt,rulecolor=\color{deitelblue!40},
  breaklines=true,showstringspaces=false,
  keywordstyle=\color{deitelblue}\bfseries,
  commentstyle=\color{deitelgreen}\itshape,
  stringstyle=\color{deitelred},
  language=C,numbers=left,numberstyle=\tiny\color{gray},
  xleftmargin=18pt,xrightmargin=5pt,stepnumber=1,
}

\pagestyle{fancy}\fancyhf{}
\fancyhead[L]{\small\color{deitelblue}\textbf{C: Como Programar -- 6ª ed. | Deitel \& Deitel}}
\fancyhead[R]{\small\color{deitelblue}Cap. 13 -- Autorrevisão}
\fancyfoot[C]{\small\thepage}
\renewcommand{\headrulewidth}{0.4pt}

\titleformat{\section}[block]{\large\bfseries\color{deitelblue}}{\thesection.}{0.5em}{}[\titlerule]
\titleformat{\subsection}[block]{\normalsize\bfseries\color{deitelblue!80}}{}{}{}

\hypersetup{colorlinks=true,linkcolor=deitelblue}

\begin{document}

\begin{titlepage}
  \centering\vspace*{2cm}
  {\color{deitelblue}\rule{\linewidth}{2pt}}\\[0.4cm]
  {\huge\bfseries\color{deitelblue}C: Como Programar}\\[0.2cm]
  {\large\color{deitelblue}6ª Edição $\cdot$ Paul Deitel \& Harvey Deitel}\\[0.2cm]
  {\color{deitelblue}\rule{\linewidth}{2pt}}\\[1cm]
  {\LARGE\bfseries Capítulo 13}\\[0.3cm]
  {\Large\color{deitelblue!80}O Pré-processador em C}\\[0.8cm]
  \begin{tcolorbox}[colback=deitellightblue,colframe=deitelblue,width=0.85\linewidth]
    \centering{\large\bfseries Exercícios de Autorrevisão (13.1--13.3)}\\[0.2cm]
    {\normalsize Resolvidos passo a passo, abordagem incremental\\
    Capítulos 1 a 13}
  \end{tcolorbox}
\end{titlepage}

\section*{Construindo sobre os Capítulos Anteriores}

No Capítulo~13, voltamos ao \texttt{\#include} e \texttt{\#define} que usamos
desde o Cap.~2 -- agora com profundidade total. O \textbf{pré-processador} é um
programa que roda \textit{antes} do compilador, modificando o código-fonte:

\begin{itemize}[noitemsep]
  \item \textbf{Diretivas} -- linhas que começam com \texttt{\#}, processadas
  \textit{antes} da compilação
  \item \textbf{Constantes simbólicas} -- \texttt{\#define} para nomes legíveis
  no lugar de valores literais
  \item \textbf{Macros} -- substituições com parâmetros, parecidas com funções
  mas expandidas em linha
  \item \textbf{Inclusão de arquivos} -- \texttt{\#include <\ldots>} (sistema) e
  \texttt{\#include "\ldots"} (usuário)
  \item \textbf{Compilação condicional} -- \texttt{\#if}, \texttt{\#ifdef},
  \texttt{\#ifndef}, \texttt{\#elif}, \texttt{\#else}, \texttt{\#endif}
  \item \textbf{Operadores de pré-processamento} -- \texttt{\#} (\textit{stringize})
  e \texttt{\#\#} (\textit{token paste})
  \item \textbf{Constantes predefinidas} -- \texttt{\_\_LINE\_\_}, \texttt{\_\_FILE\_\_},
  \texttt{\_\_DATE\_\_}, \texttt{\_\_TIME\_\_}, \texttt{\_\_STDC\_\_}
  \item \textbf{\texttt{assert}} -- macro de diagnóstico para depuração
\end{itemize}

\begin{boaprática}
  Macros são poderosas mas perigosas. Sempre que possível, prefira:
  \texttt{const} a \texttt{\#define} para constantes; \texttt{inline} (C99+) a
  macros para pequenas funções. Reserve macros para o que só elas podem fazer:
  compilação condicional, geração de código, debug.
\end{boaprática}

\tableofcontents\newpage

% ============================================================
\section{Exercício 13.1 -- Preenchimento de Lacunas (13 itens)}
% ============================================================

\begin{respostabox}
\begin{enumerate}[label=\textbf{\alph*)}]

\item Toda diretiva do pré-processador começa com \textbf{\texttt{\#}}.

\item Compilação condicional estendida usa as diretivas \textbf{\texttt{\#elif}}
e \textbf{\texttt{\#else}}.

\item A diretiva \textbf{\texttt{\#define}} cria macros e constantes simbólicas.

\item Apenas caracteres de \textbf{espaço em branco} podem aparecer antes de uma
diretiva na mesma linha.

\item A diretiva \textbf{\texttt{\#undef}} descarta nomes de constantes simbólicas e macros.

\item As notações abreviadas para \texttt{\#if defined(nome)} e
\texttt{\#if !defined(nome)} são \textbf{\texttt{\#ifdef}} e
\textbf{\texttt{\#ifndef}}.

\item A \textbf{compilação condicional} permite controlar a execução das diretivas
e a compilação do código.

\item A macro \textbf{\texttt{assert}} imprime mensagem e termina a execução se
a expressão for 0.

\item A diretiva \textbf{\texttt{\#include}} insere o conteúdo de um arquivo em outro.

\item O operador \textbf{\texttt{\#\#}} (\textit{token paste}) concatena seus dois argumentos.

\item O operador \textbf{\texttt{\#}} (\textit{stringize}) converte seu operando em string.

\item O caractere \textbf{\texttt{\textbackslash}} (barra invertida) indica que o
texto de substituição continua na próxima linha.

\item A diretiva \textbf{\texttt{\#line}} faz com que as linhas sejam renumeradas
a partir do valor indicado.

\end{enumerate}

\textbf{Resumo das diretivas:}
\begin{center}
\renewcommand{\arraystretch}{1.25}
\begin{tabular}{l l}
\toprule
\textbf{Diretiva} & \textbf{Finalidade} \\ \midrule
\texttt{\#include}   & inclui conteúdo de arquivo \\
\texttt{\#define}    & define constante ou macro \\
\texttt{\#undef}     & remove definição \\
\texttt{\#ifdef} / \texttt{\#ifndef} & teste de definição \\
\texttt{\#if} / \texttt{\#elif} / \texttt{\#else} / \texttt{\#endif} & compilação condicional \\
\texttt{\#error}     & gera erro de compilação \\
\texttt{\#pragma}    & instrução específica do compilador \\
\texttt{\#line}      & renumera linhas \\
\bottomrule
\end{tabular}
\end{center}
\end{respostabox}

\newpage

% ============================================================
\section{Exercício 13.2 -- Constantes Predefinidas}
% ============================================================

\begin{respostabox}
\begin{lstlisting}
/* Ex13.2: macros_predefinidas.c
   Imprime os valores das constantes simbolicas predefinidas */
#include <stdio.h>

int main( void )
{
    printf( "__LINE__ = %d\n", __LINE__ );
    printf( "__FILE__ = %s\n", __FILE__ );
    printf( "__DATE__ = %s\n", __DATE__ );
    printf( "__TIME__ = %s\n", __TIME__ );
    printf( "__STDC__ = %d\n", __STDC__ );

    return 0;
}
\end{lstlisting}

\textbf{Saída típica:}
\begin{saida}
\_\_LINE\_\_ = 5\\
\_\_FILE\_\_ = macros\_predefinidas.c\\
\_\_DATE\_\_ = May 13 2026\\
\_\_TIME\_\_ = 14:23:45\\
\_\_STDC\_\_ = 1
\end{saida}

\textbf{Significado de cada constante:}
\begin{itemize}[noitemsep]
  \item \texttt{\_\_LINE\_\_} -- número da linha atual no código-fonte (inteiro).
  \item \texttt{\_\_FILE\_\_} -- nome do arquivo-fonte atual (string).
  \item \texttt{\_\_DATE\_\_} -- data da compilação (string, formato ``Mmm dd yyyy'').
  \item \texttt{\_\_TIME\_\_} -- hora da compilação (string, formato ``hh:mm:ss'').
  \item \texttt{\_\_STDC\_\_} -- vale 1 se compilador adere ao padrão ANSI/ISO C.
\end{itemize}

\textbf{Aplicações práticas:}
\begin{itemize}[noitemsep]
  \item \textbf{Logs com origem:} \texttt{printf("Erro em \%s:\%d", \_\_FILE\_\_, \_\_LINE\_\_)}.
  \item \textbf{Identificação de build:} mostrar data/hora da compilação no
  \texttt{--version} do programa.
  \item \textbf{Portabilidade:} checar \texttt{\_\_STDC\_\_} antes de usar
  recursos não-portáteis.
\end{itemize}

\begin{boaprática}
  Essas macros são a base da macro \texttt{assert} e dos sistemas de log em C.
  Sempre que precisar identificar onde algo ocorreu no código, lembre-se de
  \texttt{\_\_FILE\_\_} e \texttt{\_\_LINE\_\_}.
\end{boaprática}
\end{respostabox}

\newpage

% ============================================================
\section{Exercício 13.3 -- Diretivas do Pré-processador}
% ============================================================

\begin{respostabox}
\begin{enumerate}[label=\textbf{\alph*)}]

\item \textbf{Definir \texttt{YES} como 1:}
\begin{lstlisting}
#define YES 1
\end{lstlisting}

\item \textbf{Definir \texttt{NO} como 0:}
\begin{lstlisting}
#define NO 0
\end{lstlisting}

\item \textbf{Incluir \texttt{common.h} (mesmo diretório):}
\begin{lstlisting}
#include "common.h"
\end{lstlisting}
Diferença crucial:
\begin{itemize}[noitemsep]
  \item \texttt{\#include "arquivo.h"} -- procura primeiro no diretório do
  arquivo-fonte (cabeçalhos do projeto).
  \item \texttt{\#include <arquivo.h>} -- procura nos diretórios do sistema
  (cabeçalhos da biblioteca-padrão).
\end{itemize}

\item \textbf{Renumerar a partir de 3000:}
\begin{lstlisting}
#line 3000
\end{lstlisting}
Útil para geradores de código que produzem arquivos C, para que mensagens de
erro do compilador apontem para o arquivo \textit{original} (não o gerado).

\item \textbf{Redefinir \texttt{TRUE} se estiver definida (sem \texttt{\#ifdef}):}
\begin{lstlisting}
#if defined( TRUE )
    #undef TRUE
    #define TRUE 1
#endif
\end{lstlisting}

\item \textbf{Redefinir \texttt{TRUE} se definida (com \texttt{\#ifdef}):}
\begin{lstlisting}
#ifdef TRUE
    #undef TRUE
    #define TRUE 1
#endif
\end{lstlisting}

\item \textbf{Definir \texttt{FALSE} dependendo de \texttt{TRUE}:}
\begin{lstlisting}
#if TRUE
    #define FALSE 0
#else
    #define FALSE 1
#endif
\end{lstlisting}

\item \textbf{Macro \texttt{VOLUME\_CUBO}:}
\begin{lstlisting}
#define VOLUME_CUBO( x ) ( ( x ) * ( x ) * ( x ) )
\end{lstlisting}

\begin{erroco}
  \textbf{Parênteses são essenciais em macros!} Sem eles:
\begin{lstlisting}
#define VOLUME_CUBO( x )  x * x * x  /* PERIGOSO */

VOLUME_CUBO( 1 + 2 )  /* expande para: 1 + 2 * 1 + 2 * 1 + 2 = 7 */
                      /* deveria ser:   3 * 3 * 3 = 27         */
\end{lstlisting}
Com parênteses corretos:
\begin{lstlisting}
VOLUME_CUBO( 1 + 2 )  /* expande para: (1+2)*(1+2)*(1+2) = 27 */
\end{lstlisting}
\textbf{Regra de ouro:} sempre envolver cada parâmetro e a expressão inteira
em parênteses.
\end{erroco}

\end{enumerate}
\end{respostabox}

\newpage

% ============================================================
\section*{Gabarito Rápido -- Capítulo 13: Autorrevisão}

\begin{tcolorbox}[colback=deitellightblue,colframe=deitelblue,title={\bfseries\large Respostas Resumidas}]

\textbf{13.1:}
\begin{itemize}[noitemsep]
  \item[a)] \texttt{\#} \quad b) \texttt{\#elif}, \texttt{\#else} \quad c) \texttt{\#define}
  \item[d)] espaço em branco \quad e) \texttt{\#undef} \quad f) \texttt{\#ifdef}, \texttt{\#ifndef}
  \item[g)] compilação condicional \quad h) \texttt{assert} \quad i) \texttt{\#include}
  \item[j)] \texttt{\#\#} \quad k) \texttt{\#} \quad l) \texttt{\textbackslash}
  \quad m) \texttt{\#line}
\end{itemize}

\medskip\textbf{13.2:} Programa imprime \texttt{\_\_LINE\_\_}, \texttt{\_\_FILE\_\_},
\texttt{\_\_DATE\_\_}, \texttt{\_\_TIME\_\_}, \texttt{\_\_STDC\_\_} -- todas as
constantes predefinidas pelo padrão.

\medskip\textbf{13.3 -- 8 diretivas:}
\begin{itemize}[noitemsep]
  \item[a)] \texttt{\#define YES 1}
  \item[b)] \texttt{\#define NO 0}
  \item[c)] \texttt{\#include "common.h"}
  \item[d)] \texttt{\#line 3000}
  \item[e)] \texttt{\#if defined(TRUE) ... \#undef ... \#define ... \#endif}
  \item[f)] \texttt{\#ifdef TRUE ... \#undef ... \#define ... \#endif}
  \item[g)] \texttt{\#if TRUE ... \#else ... \#endif}
  \item[h)] \texttt{\#define VOLUME\_CUBO(x) ((x)*(x)*(x))}
\end{itemize}

\end{tcolorbox}

\end{document}
